%------------------------
% Modèle de CV
% Auteur : Anubhav Singh
% Github : https://github.com/xprilion
% Licence : MIT
%------------------------

\documentclass[a4paper,20pt]{article}

\usepackage{latexsym}
\usepackage[empty]{fullpage}
\usepackage{titlesec}
\usepackage{marvosym}
\usepackage[usenames,dvipsnames]{color}
\usepackage{verbatim}
\usepackage{enumitem}
\usepackage[pdftex]{hyperref}
\usepackage{fancyhdr}

\pagestyle{fancy}
\fancyhf{} % Efface tous les champs d'en-tête et de pied de page
\fancyfoot{}
\renewcommand{\headrulewidth}{0pt}
\renewcommand{\footrulewidth}{0pt}

% Ajustement des marges
\addtolength{\oddsidemargin}{-0.530in}
\addtolength{\evensidemargin}{-0.375in}
\addtolength{\textwidth}{1in}
\addtolength{\topmargin}{-.45in}
\addtolength{\textheight}{1in}

\urlstyle{rm}

\raggedbottom
\raggedright
\setlength{\tabcolsep}{0in}

% Formatage des sections
\titleformat{\section}{
  \vspace{-10pt}\scshape\raggedright\large
}{}{0em}{}[\color{black}\titlerule \vspace{-5pt}]

%-------------------------
% Commandes personnalisées
\newcommand{\resumeItem}[2]{
  \item\small{
    \textbf{#1}{: #2 \vspace{-2pt}}
  }
}

\newcommand{\resumeItemWithoutTitle}[1]{
  \item\small{
    {\vspace{-2pt}}
  }
}

\newcommand{\resumeSubheading}[4]{
  \vspace{-1pt}\item
    \begin{tabular*}{0.97\textwidth}{l@{\extracolsep{\fill}}r}
      \textbf{#1} & #2 \\
      \textit{#3} & \textit{#4} \\
    \end{tabular*}\vspace{-5pt}
}

\newcommand{\resumeSubItem}[2]{\resumeItem{#1}{#2}\vspace{-3pt}}

\renewcommand{\labelitemii}{$\circ$}

\newcommand{\resumeSubHeadingListStart}{\begin{itemize}[leftmargin=*]}
\newcommand{\resumeSubHeadingListEnd}{\end{itemize}}
\newcommand{\resumeItemListStart}{\begin{itemize}}
\newcommand{\resumeItemListEnd}{\end{itemize}\vspace{-5pt}}

%-----------------------------
%%%%%%  LE CV COMMENCE ICI  %%%%%%

\begin{document}

%----------EN-TÊTE-----------------
\begin{tabular*}{\textwidth}{l@{\extracolsep{\fill}}r}
  \textbf{{\LARGE Mohamed Youssef Chlendi}} & Email : \href{mailto:}{youssef.chlendi@gmail.com}\\
  &
   Mobile :~~~+216 22 079 448 \\
  \href{https://github.com/youssefchlendi}{Github : ~~github.com/youssefchlendi} \\
  
  \href{https://www.linkedin.com/in/mohamed-youssef-chlendi/}{LinkedIn : ~~Mohamed Youssef Chlendi}
\end{tabular*}


	    
\vspace{-5pt}
\section{Résumé des Compétences}
\resumeSubHeadingListStart
\resumeSubItem{Langages}{~~~~~~~~~~~~PHP, JavaScript, TypeScript, Java, Python, SQL, HTML, CSS, C, C\#, Dart, Kotlin, PowerShell}
\resumeSubItem{Frameworks}{~~~~~~~~~~Laravel, Symfony, Spring Boot, ASP.NET, .NET Core, ExpressJs, Flutter, React Native, .NET MAUI, .NET 4.8, LAMP, VueJs, Angular, Ionic, Bootstrap, React}
\resumeSubItem{Outils}{~~~~~~~~~~~~~~~~~~~~GIT, MySQL, POSTMAN, GitHub, ZenHub, Azure DevOps, Firebase, Figma, Adobe Illustrator, GDPicture, WebSockets, Docker, Exchange Server}
\resumeSubItem{Bases de Données}{~~~~~~~~~~~~~~MySQL, MongoDB, SQL Server, PostgreSQL}
\resumeSubItem{Systèmes d'Exploitation}{~~Linux (Debian), Windows, macOS}
\resumeSubItem{Compétences Transversales}{~~~~~~~~~~~~~Leadership, Gestion d'événements, Gestion du temps, Analyse des besoins clients, Communication}
\resumeSubItem{Méthodologie}{~~~~~~~~Scrum, Agile, DevOps}
\resumeSubItem{Sécurité}{~~~~~~~~~~~~~~OIDC, OAuth 2.0, Authentification par formulaires, SSO}
\resumeSubHeadingListEnd

%----------EXPÉRIENCES-----------------
\vspace{-5pt}
\section{Expériences Professionnelles}
\resumeSubHeadingListStart
\vspace{-5pt}
\resumeSubheading{Capgemini}{À distance}
    {Développeur Senior .NET/Vue}{Septembre 2025 - Présent}
    
    \vspace{4pt}
    \textbf{Projet 1 : Deliverable Progressor (DP) - Gestion de Packages Électroniques intégrée à Teamcenter Siemens} \hfill \textit{Équipe de 3 développeurs}
    \resumeItemListStart
        \resumeItem{Leadership Technique}{Lead technique : mise en place des bonnes pratiques frontend, règles de tests unitaires/intégration et validation des PRs pour garantir la qualité du code.}
        \resumeItem{Migration Vue 2 vers Vue 3.3}{Direction de la migration vers Vue 3.3 : refactorisation vers la Composition API, mise à jour de l'écosystème et établissement de standards modernes pour l'équipe.}
        \resumeItem{Client API Teamcenter (Siemens)}{Conception d'un client API Teamcenter avec gestion centralisée des erreurs pour la synchronisation des MCUs, groupes de classification, hiérarchies de pièces (LIB01) et types de logiciels.}
        \resumeItem{Architecture d'Authentification}{Conception de l'architecture complète basée sur \textbf{Azure AD (Microsoft Entra ID)}, \textbf{OIDC} et \textbf{PBAC}, assurant une gestion fine des permissions par rôle et attribut.}
        \resumeItem{APIs Backend - Gestion de Packages Logiciels}
          {Développement des APIs pour la création et gestion de packages : pagination, tri et filtrage avancés multi-critères, migration des relations Item-Package vers du Many-to-Many, gestion des workflows de cycle de vie (état production, permissions de téléchargement liées à l'attribut lifetime).}
        \resumeItem{Module de Gestion des Composants Logiciels}
          {Implémentation complète du module des composants logiciels : création indépendante du package, liaison d'items existants post-création, upload automatique avec chemin de fichier configurable, journalisation détaillée et métadonnées enrichies (dates de création/modification, éditeur).}
        \resumeItem{Supervision DevOps et CI/CD}
          {Mise en place et maintenance des pipelines CI/CD, optimisation Docker, et gestion de la migration du versioning de Bitbucket vers GitHub avec mise en place des workflows de validation automatique.}
    \resumeItemListEnd
    \vspace{4pt}
    \textbf{Projet 2 : ATUM - Gestion de Binaires pour Appareils Électroménagers (Electrolux)} \hfill \textit{Équipe de 8 développeurs}
    \resumeItemListStart
        \resumeItem{Firmware et OTA}{Développement de l'UI de configuration firmware et génération automatique de binaires ; mise en place du système de distribution de mises à jour OTA.}
        \resumeItem{Refactorisation Visualiseur STL}{Remplacement de la solution existante par \textbf{Three.js} : réduction du temps de chargement de \textbf{70\%+} (50s → \textasciitilde10s) sur le HOB PowerController.}
        \resumeItem{Migration .NET 10}{Direction de la migration .NET 9 → .NET 10 avec mise à jour complète des dépendances et adaptation aux nouvelles versions de C\#.}
        \resumeItem{Infrastructure Azure}{Intégration Azure (Entra ID, Azure Web Services, Portainer) ; conduite d'une étude de migration de l'infrastructure VM Azure vers une architecture cloud native.}
        \resumeItem{Automatisation du Déploiement}{Pipeline entièrement automatisé via \textbf{Docker, Ansible et CI/CD} (déploiement manuel → quelques minutes) ; migration du versioning de Bitbucket vers GitHub avec évaluation d'Azure DevOps.}
        \resumeItem{Besoins Clients et Recrutement}{Participation directe aux réunions clients pour extraction et rédaction des specs ; conduite de \textbf{10 entretiens techniques} pour le processus de recrutement.}
    \resumeItemListEnd

    {Technologies : .NET 9/10, C\# 13, Vue.js 2/3.3, Docker, Ansible, Azure Cloud, Azure AD (Microsoft Entra ID), OIDC, PBAC, SQL Server, Three.js, Teamcenter API (Siemens), GitHub Actions, CI/CD}


\resumeSubheading{NeoLedge, Zarzis, Tunisie}{À distance}
    {Développeur Logiciel}{Mai 2023 - Septembre 2025}
    \resumeItemListStart
        \resumeItem{Système de Signature Électronique de Documents}
          {Conception et implémentation d'une solution complète de signature électronique comprenant : une CLI pour signer des fichiers PDF conformément à la norme PADES-B, des extensions web pour l'intégration aux navigateurs, et l'intégration avec des applications web pour permettre un système de signature de documents en ligne de bout en bout. Solution certifiée avec succès auprès des autorités nationales.}
        \resumeItem{Boîte à Outils Multi-Plateformes}
          {Création et publication d'une boîte à outils enveloppant la CLI pour Windows et macOS, garantissant la compatibilité avec les architectures ARM et Intel.}
        \resumeItem{Développement d'Installateurs}
          {Création d'installateurs pour la boîte à outils et ses dépendances, permettant un déploiement fluide sur Windows et macOS.}
        \resumeItem{Intégration d'Extensions Web}
          {Développement d'un système pour rendre la CLI accessible depuis des applications web via des hôtes de messagerie natives et des extensions compatibles avec les principaux navigateurs.}
        \resumeItem{Développement d'Extension Outlook}
          {Conception et implémentation d'une extension Outlook interagissant avec le serveur Exchange, gérant l'authentification OIDC et par formulaire pour optimiser les flux de gestion documentaire.}
        \resumeItem{Maintenance du Système GED d'Entreprise}
          {Participation à la maintenance et au développement d'ELISE, la solution principale de Gestion Électronique de Documents (GED) de l'entreprise, assurant la stabilité du système et l'amélioration des fonctionnalités.}
        \resumeItem{Analyse des Besoins Clients}
          {Réalisation d'études approfondies des cas d'utilisation clients et préparation de rapports détaillés pour assurer une implémentation fluide et une réponse rapide aux besoins des clients.}
        \resumeItem{Pipeline CI/CD}
          {Mise en place et maintenance de pipelines CI/CD avec Azure DevOps pour une intégration et un déploiement fluide des nouvelles fonctionnalités.}
        \resumeItem{Applications Mobiles d'Entreprise}
          {Développement et maintenance d'applications mobiles d'entreprise, intégration des services web SOAP pour une communication backend transparente.}
        \resumeItem{Certification et Conformité}
          {Facilitation de l'homologation des applications via des tests rigoureux et adaptations aux normes eIDAS, avec certification obtenue par TunTrust et ANCE Tunisie.}
        \resumeItem{Recherche en Intégration Système}
          {Réalisation de recherches approfondies pour identifier des solutions optimales pour connecter les applications web aux environnements de bureau. Les solutions évaluées comprennent :
            \begin{itemize}
              \item Communication basée sur URI pour une intégration simple et flexible.
              \item Communication HTTPS côté client pour un échange de données sécurisé et robuste.
              \item Communication WebSocket pour des interactions bidirectionnelles en temps réel.
            \end{itemize}
          }
        \resumeItem{R&D pour la Modernisation d'Applications Mobiles}
          {Conduite d'une étude pour migrer une application mobile Xamarin vers une nouvelle technologie. Évaluation des options suivantes :
            \begin{itemize}
              \item \textbf{Flutter :} Rapidité de développement, fonctionnalité de rechargement à chaud et prise en charge multi-plateforme.
              \item \textbf{React Native :} Large écosystème, forte communauté et performance via des modules natifs.
              \item \textbf{.NET MAUI :} Intégration au sein de l'écosystème .NET, compatibilité ascendante et adaptée aux solutions d'entreprise.
            \end{itemize}
          Un rapport détaillé des avantages et inconvénients a été présenté à la direction.}
        \resumeItem{Maintenance d'Applications Legacy}
          {Maintenance et amélioration d'une application UWP legacy pour assurer un support client continu et la stabilité du système.}
    \resumeItemListEnd
        {Technologies : C\#, .NET Core, Développement CLI, Développement Desktop, Développement d'Extensions, Signature Électronique, Développement Web, Vue.js, Azure DevOps, Microsoft Azure, OIDC, Exchange Server, Gestion Documentaire}
\resumeSubHeadingListEnd

%----------ÉDUCATION-----------------
\vspace{-5pt}
\section{Formation}
\resumeSubHeadingListStart
    \resumeSubheading
    {Institut Supérieur d'Informatique}{Ariana, Tunisie}
    {Master en Ingénierie Logicielle Open Source}{Sep 2023 - Juin 2025}
    \resumeSubheading
    {Institut Supérieur des Études Technologiques}{Bizerte, Tunisie}
    {Licence en Technologie des Systèmes d'Information, Mention Excellent}{Sep 2020 - Juin 2023}
\resumeSubHeadingListEnd

%----------PROJETS-----------------
\vspace{-5pt}
\section{Projets}
\resumeSubHeadingListStart
\resumeSubItem{Suivi de Vélos et Service Après-Vente}{(Oct 2022 - Jan 2023)
 Participation au développement de l'application web/mobile Pixii Motors, permettant aux propriétaires de scooters de suivre et gérer leurs motos à distance. Mise en place d'un système après-vente basé sur un système de tickets.
\\ Tech : SpringBoot, Nodejs, ExpressJs, MongoDB, Microservice, Angular, Flutter, GitHub}
\vspace{2pt}
\resumeSubItem{Système pour Scouts}{(Juin 2022 - aujourd'hui)
 Participation au développement d'un projet pour gérer les inscriptions et adhésions, les ressources monétaires et matérielles, les événements, et les groupes des utilisateurs.
\\ Tech : Nuxt, TypeScript, Vuetify, Laravel, Zenhub, GitHub}
\resumeSubHeadingListEnd

%----------CERTIFICATIONS-----------------
\vspace{-5pt}
\section{Certifications}
\resumeSubHeadingListStart
\resumeSubItem{Badge Essentials Google Cloud}
    {Google Cloud Education; Délivré : Février 2022; ID : 1740267}
    {\href{https://www.cloudskillsboost.google/public_profiles/b6cf5253-6dac-4b50-8a0d-7129daad930b/badges/1740267}{Voir la Certification}}
\resumeSubItem{Certificat Scrum Fundamentals}
    {SCRUMstudy; Délivré : Février 2022; ID : 900581}
    {\href{https://c46e136a583f7e334124-ac22991740ab4ff17e21daf2ed577041.ssl.cf1.rackcdn.com/Certificate/ScrumFundamentalsCertified-MohamedYoussefChlendi-900581.pdf}{Voir la Certification}}
\resumeSubItem{Advanced C\# Programming in .NET Core}
    {Coursera / EDUCBA; Délivré : Novembre 2025}
    {\href{https://coursera.org/verify/AHZ3PB7VHVW0}{Voir la Certification}}
\resumeSubItem{Microsoft Azure for ASP.NET Core}
    {Coursera / Packt; Délivré : Septembre 2025}
    {\href{https://coursera.org/verify/HTH0OZA64L2Z}{Voir la Certification}}
\resumeSubItem{Performance Optimization and Scalability}
    {Coursera / Microsoft; Délivré : Janvier 2026}
    {\href{https://coursera.org/verify/QW1S6J4EDX8C}{Voir la Certification}}
\resumeSubItem{AWS Artificial Intelligence Practitioner Learning Plan}
    {AWS Training \& Certification; Délivré : Septembre 2025}{}
\resumeSubItem{Aspiring Architect}
    {Capgemini University; Délivré : Janvier 2026}{}
\resumeSubHeadingListEnd

%----------EXPÉRIENCE BÉNÉVOLE-----------------
\vspace{-2pt}
\section{Expérience Bénévole}
\resumeSubHeadingListStart
\resumeItem{Vice-Président}{RAS IEEE ISET Bizerte SB}
\resumeItem{Responsable Média et Communication}{Tunisian SolidWorks Users Group ISET Bizerte}
\resumeSubHeadingListEnd

\end{document}
